\chapter*{Abstract}
\addcontentsline{toc}{chapter}{Abstract}

Bangla is spoken in many regional dialects, but automatic dialect identification
for it remains little studied. This thesis classifies 5-second audio clips into
nine classes---eight regional dialects and the standard (Formal) register---using
a 13.1-hour corpus of 9,303 clips collected from television, film and online
media. The corpus has no speaker labels, so the same speaker can appear in both
training and test data; the thesis treats this as a measurement problem rather
than a footnote.

Two routes to dialect identity are compared. The \emph{acoustic route} learns
from the sound alone: hand-crafted MFCC statistics, frozen Whisper-small encoder
embeddings, and an end-to-end fine-tuned XLS-R-300M model reach test macro-F1 of
0.727, 0.917 and 0.952 on a conventional clip-level split. A strict split built
from ECAPA-TDNN pseudo-speaker clusters, calibrated against the known sources of
the Formal clips, lowers these to 0.480, 0.793 and 0.808: 12--14 macro-F1 points
of the pretrained models' clip-level scores came from speaker and channel
overlap. Fine-tuning Whisper-medium and two Bangla-adapted versions of it
raises the strict-split score to 0.92 with a leakage drop of only 6--7 points;
the three are indistinguishable, so Bangla-specific pretraining adds nothing
measurable, while a smaller Bangla FastConformer trails at 0.75.
Channel-suppression augmentation reduces XLS-R's score (0.808 to 0.748), and
frozen Whisper embeddings trained on 10\,\% of the labels outperform MFCC
features trained on all of them.

The \emph{lexical route} transcribes every clip with five off-the-shelf Bangla
speech recognisers and classifies the words. Measured against human
transcripts of dialect speech, the best recogniser---one fine-tuned on regional
speech by its publishers---reproduces 27\,\% of dialect words; recognisers
trained on standard Bangla reproduce 1--4\,\%, rewriting dialect into the
standard form. Matching regional word lists therefore identifies the dialect
poorly (macro-F1 0.36 under the strict split), and a text classifier over the
full transcript reaches 0.63. Combined with the XLS-R model, however, the
words raise the strict-split score from 0.808 to 0.856, and the clip-level score
from 0.952 to 0.967; the clip-level gain is significant, the strict-split gain
borderline.

On clips this short, Bangla dialect identity lies mostly beyond the words: the
acoustic route is the stronger single route, and the lexical route is limited by
speech recognition rather than by classification, yet the two are complementary.
All code, splits and per-clip predictions are released for reproduction.

\vspace{1em}
\noindent\textbf{Keywords:} dialect identification, Bangla, self-supervised
speech models, data leakage, speech recognition, low-resource languages
